\section*{अध्याय \ref{ch:b3:spectral} --- संहत संकारक और
स्पेक्ट्रमी प्रमेय}

\begin{solution}{exo:b3:spectral:1}
(क) $T(B)$ परिमित-विमीय $\operatorname{im}T$ का कोई परिबद्ध उपसमुच्चय है: अतः
हाइने--बोरेल (\cref{cor:b3:topology:heineborel}, $\R^n$ के साथ किसी रैखिक \omterm{def:b3:topology:continuity}{समस्थितिकता} से अंतरित) से
सापेक्षतः \omterm{def:b3:spectral:compact}{संहत}।
(ख) $I$ का \omterm{def:b3:spectral:compact}{संहत} होना यह कहता है कि संवृत इकाई गोला \omterm{def:b3:spectral:compact}{संहत} है, जो रीस
की प्रमेय (दूसरा वर्ष) से ठीक परिमित विमा में होता है। यदि किसी \omterm{def:b3:spectral:compact}{संहत}
$T$ का परिबद्ध प्रतिलोम $T^{-1}$ होता, तो $I = T^{-1}T$ \omterm{def:b3:spectral:compact}{संहत} होता (\cref{prop:b3:spectral:ideal}): जो अनंत
विमा में असंभव है।
\end{solution}

\begin{solution}{exo:b3:spectral:2}
(क) $\norm{Tx}^2 = \sum\abs{\lambda_n}^2\abs{x_n}^2 \leq
\sup\abs{\lambda_n}^2\norm x^2$, जिसमें उच्चतम को साकार करने वाले $e_n$ पर निकट-बराबरी।
(ख) ($\Leftarrow$) कर्तन $T_N$ ($n \leq N$ रखिए, आगे शून्य) परिमित कोटि के हैं और
$\vertiii{T - T_N} =
\sup_{n>N}\abs{\lambda_n} \to 0$: अतः \cref{prop:b3:spectral:ideal} से \omterm{def:b3:spectral:compact}{संहत}। ($\Rightarrow$) यदि किसी उपानुक्रम के अनुदिश $\abs{\lambda_{n_k}} \geq \delta > 0$ हो:
$\norm{Te_{n_k} - Te_{n_l}}^2 = \abs{\lambda_{n_k}}^2 +
\abs{\lambda_{n_l}}^2 \geq 2\delta^2$: अतः $(Te_{n_k})$ का कोई अभिसारी उपानुक्रम नहीं।
(ग) $(\bar\lambda_n)$ के साथ विकर्ण $T^* ={}$: \omterm{def:b3:spectral:selfadjoint}{स्वसंलग्न} तभी और केवल तभी जब सभी $\lambda_n \in \R$।
तब मानक आधार $(e_n)$ अभिलक्षणिक सदिशों का कोई लांबिक-प्रसामान्य आधार
है, जिसके अभिलक्षणिक मान $\lambda_n \to
0$ हैं: अर्थात् स्पेक्ट्रमी प्रमेय अक्षरशः।
\end{solution}

\begin{solution}{exo:b3:spectral:3}
मान लीजिए $(x_n)$ परिबद्ध है। तब $(Bx_n)$ परिबद्ध है ($\vertiii B < \infty$); $S$ की संहतता
$SBx_{n_k} \to y$ निकाल लेती है; और $A$ की \omterm{def:b3:topology:continuity}{सांतत्य} $ASBx_{n_k} \to
Ay$ दे देती है: अतः $ASB$
\omterm{def:b3:spectral:compact}{संहत} है। यदि \omterm{def:b3:spectral:compact}{संहत} $T$ और $\dim H = \infty$ के साथ $ST = TS = I$ हो: तो $I = ST$ \omterm{def:b3:spectral:compact}{संहत} होता,
जो \cref{exo:b3:spectral:1}(ख) का खंडन करता है --- और वह वही कथन है, दूसरी ओर से देखा
गया।
\end{solution}

\begin{solution}{exo:b3:spectral:4}
(क) $y$ में कोशी--श्वार्ज़ से: $\abs{T_kf(x)}^2 \leq
\bigl(\int\abs{k(x,y)}^2\dd y\bigr)\norm f_2^2$; $x$ में समाकलन कीजिए (टोनेली):
$\norm{T_kf}_2 \leq \norm k_{L^2(\square)}
\norm f_2$।

(ख) कुल $e_{mn}(x,y) = e_m(x)\overline{e_n(y)}$ $L^2(\intcc01^2)$ में लांबिक-प्रसामान्य है (टोनेली द्विक समाकल को
अलग कर देती है)। \omterm{def:b3:complete:complete}{पूर्ण}: यदि सभी $e_{mn}$ के लिए $h \perp$ हो, तो प्रत्येक
$m$ के लिए फलन $y \mapsto \int h(x,y)\overline{e_m(x)}\dd
x$ (कोशी--श्वार्ज़ और टोनेली से $L^2$ में) प्रत्येक
$\overline{e_n}$ पर लांबिक है --- और जब भी $(e_n)$ कोई \omterm{def:b3:hilbert:onb}{हिल्बर्ट आधार} हो तब संयुग्म
$(\overline{e_n})$ भी होते हैं (संयुग्मन $L^2$ का कोई समदूरीय एकैकी आच्छादन है जो
लांबिकता और \omterm{def:b3:complete:complete}{पूर्णता} को सुरक्षित रखता है) --- अतः वह लगभग सर्वत्र
$0$ है; तब लगभग प्रत्येक $y$ के लिए प्रत्येक $e_m$ के सापेक्ष $h(\cdot, y) \perp$:
लगभग सर्वत्र $h(\cdot, y) = 0$: अतः $h = 0$ (टोनेली)। इसलिए $(e_{mn})$ कोई \omterm{def:b3:hilbert:onb}{हिल्बर्ट आधार}
है; $k = \sum c_{mn}e_{mn}$ का प्रसार कीजिए। कर्तन $k_N$ (सूचकांक $\leq N$) परिमित कोटि का
$T_{k_N}$ देता है (परास $\operatorname{Vect}(e_1, \dots,
e_N)$ में), और (क) से
\[
\vertiii{T_k - T_{k_N}} \leq \norm{k - k_N}_{L^2} \to 0 :
\]
अतः $T_k$ परिमित-कोटि संकारकों की कोई मानक-सीमा है: \omterm{def:b3:spectral:compact}{संहत} (\cref{prop:b3:spectral:ideal})।
\end{solution}

\begin{solution}{exo:b3:spectral:5}
(क) किसी अभिलक्षणिक सदिश पर $\lambda\norm x^2 = \langle x, Tx\rangle \geq 0$। सूत्र सभी मान $\langle x,
Tx\rangle \geq 0$ के साथ \cref{prop:b3:spectral:sanorm}
है: अतः निरपेक्ष मान अनावश्यक है।
(ख) $(x, y) \mapsto \langle x, Ty\rangle$ कोई एर्मीती धनात्मक (संभवतः अपह्रासी) अर्ध-द्विरैखिक रूप
है; और कोशी--श्वार्ज़ की सामान्य उपपत्ति ($\langle x + t\eu^{\iu\theta}y,
T(x + t\eu^{\iu\theta}y)\rangle \geq 0$ का प्रसार कीजिए और
विविक्तकर लीजिए) निश्चितता का उपयोग कभी नहीं करती।
\end{solution}

\begin{solution}{exo:b3:spectral:6}
(क) $\norm{Mf}_2 \leq \norm f_2$, और $f_n = \sqrt n\,
\mathbf 1_{\intcc{1 - 1/n}1}$ (इकाई सदिश) पर $\norm{Mf_n}_2
\geq 1 - \frac1n$: $\vertiii M = 1$; और \omterm{def:b3:spectral:selfadjoint}{स्वसंलग्न}, क्योंकि
गुणक वास्तविक है। अभिलक्षणिक मान: लगभग सर्वत्र $xf(x) = \lambda f(x)$ से अकिंचन समुच्चय
$\{x = \lambda\}$ के बाहर लगभग सर्वत्र $f = 0$ अनिवार्य हो जाता है: $L^2$ में $f =
0$।
(ख) उन्हीं $f_n$ के साथ: $\norm{Mf_n - f_n}_2 \leq \frac1n \to
0$, जबकि $f_n \rightharpoonup 0$ (स्थिर $g \in L^2$ के लिए प्रभावी
अभिसरण प्रमेय से $\abs{\langle g, f_n\rangle} \leq \norm{g\,\mathbf
1_{\intcc{1-1/n}1}}_2 \to 0$)। यदि मानक में $Mf_{n_k} \to h$ हो, तो $f_{n_k} \to h$, जिससे $h = 0$
अनिवार्य हो जाता है (दुर्बल सीमा) फिर भी $\norm h = 1$: अतः कोई अभिसारी
उपानुक्रम नहीं।
(ग) \cref{lem:b3:spectral:existence} में ठीक ``$Tx_{n_k} \to y$'' वाला निष्कर्षण ही संहतता का उपयोग करता
है; और $M$ के लिए अधिकतमकारी अनुक्रम $x = 1$ के निकट संकेंद्रित हो
जाते हैं और उनके प्रतिबिंब $0$ पर दुर्बल रूप से अभिसरित होते हैं,
मानक में कभी नहीं: संख्यात्मक परास के शिखर पर अभिलक्षणिक सदिश बस
विद्यमान ही नहीं होता।
\end{solution}

\begin{solution}{exo:b3:spectral:7}
(क) $k(x, y) = \mathbf 1_{y < x} \in
L^2(\intcc01^2)$ के साथ $V = T_k$: अतः \omterm{def:b3:spectral:compact}{संहत} (\cref{exo:b3:spectral:4})। उसका संलग्न अष्टि $\overline{k(y,
x)} = \mathbf 1_{y > x}$
वाला अष्टि संकारक है: $V^*f(x) = \int_x^1f$; $V \neq V^*$ ($f = \mathbf 1$ पर परीक्षा लीजिए)।
(ख) यदि $Vf = \lambda f$, $\lambda \ne 0$ हो: तो $Vf$ $\intcc01$ पर \omterm{def:b3:topology:continuity}{संतत} है ($\int_0^x f$ में प्रभावी
अभिसरण), अतः $f
= \frac1\lambda Vf$ का कोई \omterm{def:b3:topology:continuity}{संतत} प्रतिनिधि है; तब $Vf$ $\mathcal C^1$ है (\omterm{def:b3:topology:continuity}{संतत}
समाकल्यों के लिए कलन की मूल प्रमेय), अतः $f$ $\mathcal C^1$ है, और $f(0) = \frac1\lambda Vf(0) = 0$
के साथ $\lambda f' = f$: $C = f(0) = 0$ के साथ $f =
C\eu^{x/\lambda}$।
(ग) $V$ \omterm{def:b3:spectral:compact}{संहत} है और संभवतः $0$ को छोड़कर उसका कोई भी अभिलक्षणिक
मान नहीं है ($Vf = 0$ समाकल का अवकलन करने पर लगभग सर्वत्र $f = 0$ अनिवार्य
कर देता है --- अतः $0$ भी नहीं): स्पेक्ट्रमी मशीनरी को सचमुच
\omterm{def:b3:spectral:selfadjoint}{स्वसंलग्नता} चाहिए, केवल संहतता नहीं।
\end{solution}

\begin{solution}{exo:b3:spectral:8}
$x = \sum_ic_ie_i + z$, $z \in \ker T$ लिखिए, अतः $\langle x,
Tx\rangle = \sum_i\mu_i\abs{c_i}^2$। \emph{निचला परिबंध}: $V_k = \operatorname{Vect}(e_1, \dots,
e_k)$ के इकाई
गोलक पर $\langle x, Tx\rangle = \sum_{i\leq
k}\mu_i\abs{c_i}^2 \geq \mu_k$: अतः $V$ पर निम्निष्ठ का उच्चिष्ठ $\geq \mu_k$ है। \emph{ऊपरी
परिबंध}: मान लीजिए $\dim V = k$ और $W_k
= \overline{\operatorname{Vect}}(e_k, e_{k+1}, \dots) +
\ker T$, जिसकी सह-विमा $k - 1$ है (उसका
\omterm{thm:b3:hilbert:decomposition}{लांबिक पूरक} $V_{k-1}$ है); $V \cap W_k \neq \{0\}$ (कोटि $\leq k - 1$ के किसी रैखिक प्रतिचित्रण
$V \to
H/W_k \cong V_{k-1}$ की अष्टि अतुच्छ होती है), और किसी इकाई $x \in V\cap W_k$ के लिए $\langle x,
Tx\rangle = \sum_{i \geq k}\mu_i\abs{c_i}^2 \leq \mu_k$: अतः
$V$ पर निम्निष्ठ $\leq \mu_k$ है। मिलाकर: पहला सूत्र; और दूसरा सममित रूप
से सिद्ध होता है ($W = W_k$ पर परीक्षा लीजिए; सह-विमा $k-1$ वाले स्वेच्छ
$W$ के लिए $W \cap V_k \neq 0$ $\langle x, Tx\rangle \geq \mu_k$ वाला कोई इकाई सदिश दे देता है)। एकदिष्टता:
बिंदुशः $\langle x, Tx\rangle \leq \langle x,
Sx\rangle$ $\max\min$ के माध्यम से अंतरित हो जाता है।
\end{solution}

\begin{solution}{exo:b3:spectral:9}
$f - \lambda Tf = g$ फिर से लिखिए। $\lambda = 0$ के लिए: $f = g$, जो सदैव अद्वितीय रूप से
हल-योग्य है। $\lambda \neq 0$ के लिए: यह $(T - \frac1\lambda)f = -\frac g\lambda$ है, और $T$ (\cref{ex:b3:spectral:minxy}) के अभिलक्षणिक
मानों $\lambda_n = \bigl((n + \frac12)\pi\bigr)^{-2}$ के साथ \omterm{thm:b3:spectral:fredholm}{फ्रेडहोम विकल्प} (\cref{thm:b3:spectral:fredholm}) से: सभी $g$ के लिए
अद्वितीय हल-योग्यता तभी और केवल तभी जब प्रत्येक $n$ के लिए $\frac1\lambda \neq \lambda_n$,
अर्थात्
\[
\lambda \neq \Bigl(n + \tfrac12\Bigr)^2\pi^2
\qquad (n = 0, 1, 2, \dots).
\]
किसी अपवादिक $\lambda = (n+\frac12)^2\pi^2$ पर: हल विद्यमान हैं तभी और केवल तभी जब $g \perp \sin\bigl((n{+}\frac12)\pi x\bigr)$, और
तब वे उस ज्या के गुणज जोड़ने तक अद्वितीय हैं।
\end{solution}

\begin{solution}{exo:b3:spectral:10}
(क) $Sx = \lambda x$: निर्देशांकों की तुलना करने पर $0 = \lambda
x_1$ और $x_n = \lambda x_{n+1}$; यदि $\lambda \ne 0$
हो तो $x_1 = 0$ और आगमन से $x = 0$; और यदि $\lambda = 0$ हो, तो $Sx = 0$ $x = 0$ अनिवार्य
कर देता है ($S$ समदूरीय)। अतः कोई अभिलक्षणिक मान नहीं। $S^*x =
\lambda x$ $x_{n+1} = \lambda x_n$
के रूप में पढ़ा जाता है: $x =
x_1(1, \lambda, \lambda^2, \dots)$, जो $\ell^2$ में ठीक तब है जब $\abs\lambda < 1$:
अर्थात् अभिलक्षणिक मानों की एक पूरी विवृत चक्र।
(ख) $\norm{Se_n - Se_m} = \norm{e_{n+1} - e_{m+1}} = \sqrt2$: परिबद्ध $(e_n)$ के प्रतिबिंब का कोई कोशी उपानुक्रम नहीं;
और इसी प्रकार $S^*e_{n+1} = e_n$। दोनों में से कोई \omterm{def:b3:spectral:compact}{संहत} नहीं है।
(ग) \omterm{def:b3:spectral:compact}{संहत} \omterm{def:b3:spectral:selfadjoint}{स्वसंलग्न} संकारकों के लिए अभिलक्षणिक मान सब कुछ पकड़ लेते
हैं (\cref{thm:b3:spectral:spectral}); पर दोनों में से कोई भी परिकल्पना हटाने पर अभिलक्षणिक
मान पूरी तरह लुप्त हो सकते हैं ($S$, वोल्तेरा) या कोई चक्र भर
सकते हैं ($S^*$): सुदृढ़ वस्तु स्पेक्ट्रम $\{\lambda : T - \lambda I \text{ व्युत्क्रमणीय नहीं}\}$ है, जिसका सिद्धांत
किसी बाद के पाठ्यक्रम का विषय है।
\end{solution}

\begin{solution}{exo:b3:spectral:11}
(क) $S$ गुणांकों $\sqrt{\mu_n} \to 0$ वाला विकर्ण संकारक है: अतः \omterm{def:b3:spectral:compact}{संहत} (\cref{exo:b3:spectral:2}(ख),
$\ker T$ से पूरित आधार $(e_n)$ पर अंतरित, जहाँ $S = 0$), \omterm{def:b3:spectral:selfadjoint}{स्वसंलग्न} (वास्तविक
विकर्ण), धनात्मक ($\langle x, Sx\rangle =
\sum\sqrt{\mu_n}\abs{\langle e_n, x\rangle}^2$), और पदशः $S^2 =
T$।

(ख) $R$ $T = R^2$ के साथ क्रमविनिमेय है; और अभिलक्षणिक मान $\mu$ वाले
$T$ के किसी अभिलक्षणिक सदिश $x$ के लिए $T(Rx) = RTx = \mu Rx$, अतः परिमित-विमीय
अभिलक्षणिक समष्टि $E_\mu = \ker(T - \mu)$ $R$-स्थायी है। $E_\mu$ पर $R$ $R^2
= \mu\,\mathrm{id}$ के
साथ सममित धनात्मक है: उसके अभिलक्षणिक मान $\rho$ $\rho^2
= \mu$, $\rho \geq 0$ संतुष्ट
करते हैं: अतः सभी $\sqrt\mu$ के बराबर, और एक ही अभिलक्षणिक मान वाला कोई
विकर्णीय संकारक अदिश होता है: $E_\mu$ पर $R = \sqrt\mu\,\mathrm{id}$। $\ker T$ पर: $\norm
{Rx}^2 = \langle x, R^2x\rangle = \langle x, Tx\rangle = 0$। अतः
$R$ $\ker T$ पर और प्रत्येक अभिलक्षणिक समष्टि पर $S$ से सहमत है,
जिनका संवृत विस्तार $H$ है (स्पेक्ट्रमी प्रमेय): $R = S$।

(ग) $\sqrt G$ $e_n =
\sqrt2\sin(n\pi x)$ पर $\frac1{n\pi}$ की भाँति क्रिया करता है: वह
\[
k(x, y) = \sum_{n\geq1}\frac{2\sin(n\pi x)\sin(n\pi
y)}{n\pi},
\]
अष्टि वाला अष्टि संकारक है, जिसकी श्रेणी $L^2(\intcc01^2)$ में अभिसरित होती है
(गुणांक $\frac1{n\pi} \in \ell^2$; आंशिक-योग अष्टियाँ परिमित-कोटि आसन्नन दे देती हैं)।
किसी प्रारंभिक संवृत रूप की आवश्यकता नहीं: स्पेक्ट्रमी पक्ष \emph{ही}
वह संकारक है।
\end{solution}

\begin{solution}{exo:b3:spectral:12}
(क) $T^*T$ \omterm{def:b3:spectral:compact}{संहत} (किसी परिबद्ध और किसी \omterm{def:b3:spectral:compact}{संहत संकारक} का गुणनफल, \cref{exo:b3:spectral:3}),
\omterm{def:b3:spectral:selfadjoint}{स्वसंलग्न} ($(T^*T)^* = T^*T$) और धनात्मक ($\langle x, T^*Tx\rangle =
\norm{Tx}^2$) है। अष्टि: $T^*Tx = 0 \Rightarrow \norm{Tx}^2 =
\langle x, T^*Tx\rangle = 0 \Rightarrow Tx = 0$, और विलोमतः
$\ker T^*T = \ker T$। स्पेक्ट्रमी प्रमेय $T^*Te_n = s_n^2e_n$, $s_n > 0$ वाले लांबिक-प्रसामान्य $(e_n)$
दे देती है, जो $(\ker T)^\perp$ को फैलाते हैं।

(ख) $\langle f_m, f_n\rangle = \frac{\langle Te_m,
Te_n\rangle}{s_ms_n} = \frac{\langle e_m,
T^*Te_n\rangle}{s_ms_n} = \frac{s_n^2}{s_ms_n}\delta_{mn} =
\delta_{mn}$। $x_0 \in \ker T$ के साथ $x = x_0 + \sum_n\langle e_n, x\rangle
e_n$ का प्रसार कीजिए (संवृत विस्तार में
\omterm{thm:b3:hilbert:parseval}{पार्सेवाल} और साथ में अष्टि); \omterm{def:b3:topology:continuity}{संतत} $T$ लगाने पर:
\[
Tx = \sum_n\langle e_n, x\rangle\,Te_n =
\sum_ns_n\langle e_n, x\rangle\,f_n,
\]
जिसमें श्रेणी इसलिए अभिसरित होती है कि उसके आंशिक योग कोशी हैं ($\norm{\sum_{N<n\leq M}s_n\langle e_n, x\rangle f_n}^2 =
\sum s_n^2\abs{\langle e_n, x\rangle}^2$,
जिस पर $\sup_{n>N}s_n^2\cdot\norm x^2$ का प्रभुत्व है, और $s_n \to 0$)।

(ग) $\norm{Tx}^2 = \sum_ns_n^2\abs{\langle e_n, x\rangle}^2
\leq (\max s_n)^2\norm x^2$, जो अधिकतमकारी $e_n$ पर प्राप्त होता है: $\vertiii T = \max s_n$। एकल मान
विघटन को कोटि $N$ पर कर्तित करने से मानक $\sup_{n>N}s_n \to 0$ का संकारक बचता है:
अर्थात् परिमित-कोटि आसन्नन। वोल्तेरा संकारक का कोई अभिलक्षणिक मान
नहीं है (\cref{exo:b3:spectral:7}), पर $V^*V$ का है: $V^*Vf(x) = \int_x^1\int_0^tf(s)\,\dd s\,\dd t$ कोई सममित धनात्मक अष्टि संकारक
है (अष्टि $1 - \max(x,y)$, अर्थात् तार-प्रकार की कोई ग्रीन अष्टि), जिसके
अभिलक्षणिक युग्म --- परिसीमा मान समस्या $-u'' = \lambda^{-1}u$, $u'(0)
= u(1) = 0$, अर्थात् \cref{exo:b3:spectral:9}
के कुल के माध्यम से परिकलित --- एकल मान $s_n = \bigl((n + \frac12)\pi\bigr)^{-1}$ दे देते हैं। एकल मान
विघटन ठीक इसलिए \emph{दो} लांबिक-प्रसामान्य कुलों पर रहता है कि $V$
अपनी अभिलक्षणिक ज्यामिति को घुमाकर हटा देता है: कोई अभिलक्षणिक सदिश
नहीं, फिर भी दो भिन्न आधारों के बीच \omterm{def:b3:complete:complete}{पूर्ण} विकर्ण संरचना।
\end{solution}

\begin{solution}{pb:b3:spectral:1}
\textbf{1.} \omterm{def:b3:topology:continuity}{सांतत्य}: $\min$ और $\max$ \omterm{def:b3:topology:continuity}{संतत} हैं; सममिति: $x, y$ की
अदला-बदली न $\min(x,y)$ बदलती है न $1 -
\max(x,y)$। परिबंध: $0 \leq g$, और $g(x,y) \leq
\max\cdot(1-\max)$-प्रकार
के परिबंध $g \leq \frac14$ दे देते हैं ($u
= \max$ के लिए: $\min \leq u$ अतः $g \leq u(1-u) \leq \frac14$)। $g
\in L^2(\square)$:
हिल्बर्ट--श्मिट, अतः $G$ \omterm{def:b3:spectral:compact}{संहत} (\cref{exo:b3:spectral:4}); और अष्टि वास्तविक सममित
है: $G$ \omterm{def:b3:spectral:selfadjoint}{स्वसंलग्न}।

\textbf{2.} $y = x$ पर विभाजित करने पर:
\[
u(x) = (1 - x)\int_0^x y\,f(y)\,\dd y +
x\int_x^1(1 - y)\,f(y)\,\dd y .
\]
\omterm{def:b3:topology:continuity}{संतत} $f$ के लिए अवकलन कीजिए (गुणनफल और मूल प्रमेय):
\[
u'(x) = -\int_0^xyf + \int_x^1(1-y)f
\qquad\text{(परिसीमा पद निरस्त हो जाते हैं)},
\]
और $u''(x) = -xf(x) - (1 - x)f(x) = -f(x)$; तथा स्पष्टतः $u(0) =
u(1) = 0$। विलोमतः यदि $u \in \mathcal C^2$ दोनों सिरों पर लुप्त
हो, तो $w = u - G(-u'')$ $w'' = 0$, $w(0) = w(1) =
0$ संतुष्ट करता है: अतः $w$ आफ़ीन है और दो
बार लुप्त होता है, इसलिए $w = 0$।

\textbf{3.} मान लीजिए $Gf = 0$, $f \in L^2$। $\psi \in \mathcal
C_c^\infty(\intoo01)$ के लिए: प्रश्न 2 से $\psi = G(-\psi'')$,
अतः
\[
\langle f, \psi\rangle = \langle f, G(-\psi'')\rangle
= \langle Gf, -\psi''\rangle = 0
\]
($G$ \omterm{def:b3:spectral:selfadjoint}{स्वसंलग्न}): अतः मूल प्रमेयिका (\cref{cor:b3:lp:fundlemma}) से लगभग सर्वत्र $f = 0$।
धनात्मकता: \omterm{def:b3:topology:continuity}{संतत} $f$ के लिए $u = Gf$ के साथ
\[
\langle f, Gf\rangle = \int_0^1 fu = \int_0^1(-u'')u
= \bigl[-u'u\bigr]_0^1 + \int_0^1(u')^2 = \int_0^1(u')^2
\geq 0 ;
\]
और $f \in L^2$ के लिए $L^2$ में \omterm{def:b3:topology:continuity}{संतत} $f_n$ से आसन्नन कीजिए: दोनों पक्ष
सीमा तक चले जाते हैं ($G$ परिबद्ध)।

\textbf{4.} यदि $Gu = \lambda u$, $\lambda \neq 0$ हो: $Gu$ \omterm{def:b3:topology:continuity}{संतत} है ($\abs{Gu(x) - Gu(x')} \leq \norm{g(x,\cdot) -
g(x',\cdot)}_2\norm u_2$, और अष्टि
एकसमान रूप से \omterm{def:b3:topology:continuity}{संतत} है), अतः $u$ का कोई \omterm{def:b3:topology:continuity}{संतत} प्रतिनिधि है; तब प्रश्न
2 के सूत्र $Gu \in \mathcal C^2$ दिखा देते हैं, अतः $-\lambda u'' =
-(Gu)'' = u$ और $u(0) = u(1) = 0$ के साथ $u =
\frac1\lambda Gu \in \mathcal C^2$।

\textbf{5.} $-\lambda u'' = u$, $u(0) = 0$: $u = A\sin(x/\sqrt
\lambda)$ (धनात्मक $\lambda$: प्रश्न 3 से अभिलक्षणिक
सदिशों पर $\lambda =
\langle u, Gu\rangle/\norm u^2 > 0$)। $u(1) =
0$ $\frac1{\sqrt\lambda} = n\pi$ अनिवार्य कर देता है: $\lambda_n =
\frac1{n^2\pi^2}$, अभिलक्षणिक
फलन $\sin(n\pi x)$, मानकीकृत $e_n = \sqrt2\sin(n\pi x)$ ($\int_0^12\sin^2(n\pi x)\dd x =
1$)। लांबिकता जाँच: $m
\neq n$ के लिए $2\sin(m\pi x)\sin(n\pi x) =
\cos((m-n)\pi x) - \cos((m+n)\pi x)$
का समाकल $0$ है।

\textbf{6.} $\ker G = \{0\}$ (प्रश्न 3), अतः \cref{thm:b3:spectral:spectral}(1) $H =
\overline{\operatorname{Vect}}(e_n)$ देती है: इसलिए ज्या
$L^2(\intcc01)$ का कोई \omterm{def:b3:hilbert:onb}{हिल्बर्ट आधार} हैं। $f(x) = x(1 - x)$ के लिए:
\[
c_n = \sqrt2\int_0^1x(1-x)\sin(n\pi x)\,\dd x
= \sqrt2\;\frac{2\bigl(1 - (-1)^n\bigr)}{n^3\pi^3}
= \begin{cases}\dfrac{4\sqrt2}{n^3\pi^3} & n \text{ विषम},\\
0 & n \text{ सम},\end{cases}
\]
(दो बार खंडशः समाकलन)। \omterm{thm:b3:hilbert:parseval}{पार्सेवाल}: $\int_0^1x^2(1-x)^2\dd
x = \frac1{30} = \sum_{n \text{ विषम}}\frac{32}{n^6\pi^6}$, अर्थात् $\sum_{n\text{ विषम}}n^{-6} = \frac{\pi^6}{960}$।

\textbf{7.} $\langle e_n, Ge_n\rangle = \lambda_n\norm{e_n}^2
= \lambda_n$। स्थिर $x$ के लिए $g(x,
\cdot)$ के गुणांक: $\langle e_n, g(x,\cdot)\rangle = (Ge_n)(x) =
\lambda_ne_n(x)$, अतः
$L^2$ में $g(x,\cdot) =
\sum_n\lambda_ne_n(x)\,e_n$। स्पष्ट श्रेणी $\sum_n\lambda_ne_n(x)e_n(y) = \sum_n\frac{2\sin(n\pi
x)\sin(n\pi y)}{n^2\pi^2}$ वर्ग पर प्रसामान्य रूप से अभिसरित
होती है ($\abs{\text{पद}} \leq \frac2{n^2\pi^2}$): अतः उसका योग \omterm{def:b3:topology:continuity}{संतत} है, और प्रत्येक $x$ के लिए उसके
$L^2(\dd y)$-गुणांक $g(x, \cdot)$ के ही हैं: इसलिए दोनों \omterm{def:b3:topology:continuity}{संतत} फलन प्रत्येक $(x, y)$
के लिए सहमत हैं। $y = x$ रखकर समाकलन करने पर (प्रसामान्य अभिसरण पदशः
समाकलन की अनुमति देता है):
\[
\int_0^1g(x,x)\,\dd x =
\sum_n\lambda_n\int_0^1e_n(x)^2\dd x = \sum_n\lambda_n .
\]

\textbf{8.} $\int_0^1g(x,x)\dd x = \int_0^1x(1 - x)\dd x =
\frac16$, अतः $\sum_{n\geq1}\frac1{n^2\pi^2} = \frac16$:
\[
\zeta(2) = \sum_{n\geq1}\frac1{n^2} = \frac{\pi^2}6 .
\]

\textbf{9.} $f = \mathbf 1$ के लिए: $c_n =
\sqrt2\int_0^1\sin(n\pi x)\dd x = \sqrt2\,\frac{1 -
(-1)^n}{n\pi}$: विषम $n$ के लिए $c_n = \frac{2\sqrt2}{n\pi}$, और सम के
लिए $0$। \omterm{thm:b3:hilbert:parseval}{पार्सेवाल}: $1 = \sum_{n\text{
विषम}}\frac{8}{n^2\pi^2}$, अतः $\sum_{n\text{ विषम}}n^{-2} =
\frac{\pi^2}8$, और $\zeta(2) = \frac{\pi^2}8\cdot\frac{1}{1
- \frac14} = \frac{\pi^2}6$ (सम पद $\frac14\zeta(2)$ हैं)।
तुलना: किसी एक $f$ के लिए \omterm{thm:b3:hilbert:parseval}{पार्सेवाल} $\abs{\langle e_n, f\rangle}^2$ का योग लेता है; और संचिह्न
सूत्र \emph{अष्टि} के विकर्ण का समाकलन करता है, जो एक ही बार में
पूरे लांबिक-प्रसामान्य कुल पर \omterm{thm:b3:hilbert:parseval}{पार्सेवाल} का योग लेने के बराबर है ---
$\sum_n\langle e_n, Ge_n\rangle$ --- और इसलिए वह परीक्षण फलन के किसी विशेष चयन के प्रति अंधा
है।

\textbf{10.} $\sum\abs{c_n}^2 < \infty$ के साथ $\abs{c_n\cos(n\pi t)} \leq \abs{c_n}$: प्रत्येक $t$ के लिए श्रेणी $L^2$
में अभिसरित होती है (लांबिक-प्रसामान्य प्रसार, \cref{thm:b3:hilbert:parseval}(3)); और पुच्छ
परिबंध $\norm{u(t) - u_N(t)}_2^2 \leq \sum_{n>N}\abs{c_n}^2$ $t$ में एकसमान है, तथा प्रत्येक आंशिक योग $t$ में
\omterm{def:b3:topology:continuity}{संतत} है (परिमित संख्या के कोज्या): अतः $t \mapsto u(t,\cdot)$ $L^2$ में \omterm{def:b3:topology:continuity}{संतत} है। $f = \sum_{n\leq N}c_ne_n$
के लिए: प्रत्येक विधा $\cos(n\pi t)\sin(n\pi x)$ अपने पर लगाई गई $\partial_t^2 =
-n^2\pi^2 = \partial_x^2$ संतुष्ट करती है,
$x = 0,
1$ पर लुप्त होती है, और $t =
0$ पर उसका मान $\sin(n\pi x)$ तथा समय-अवकलज
$0$ है: अतः परिमित योग सब कुछ हल कर देता है। संगीत की दृष्टि से:
तार की गति \emph{स्थायी तरंगों} $e_n$ का अध्यारोपण है, जिनकी आवृत्तियाँ
$n\pi$ मूल स्वर और उसके अधिस्वर हैं; और स्पेक्ट्रमी प्रमेय कहती है
कि प्रत्येक प्रारंभिक आकृति इन शुद्ध स्वरों में अद्वितीय रूप से विघटित
हो जाती है, जिसमें गुणांक $c_n$ ही स्वर-गुण हैं। तार सुनना कोई
लांबिक-प्रसामान्य प्रसार परिकलित करना है।

\textbf{11.} $a_n = \abs{\langle u_n, x\rangle}^2$ के साथ: $m_{p+1} = \sum_n\mu_n^{p+1}a_n = \sum_n\bigl(\mu_n^{p/2}
\sqrt{a_n}\bigr)\bigl(\mu_n^{p/2+1}\sqrt{a_n}\bigr) \leq
\sqrt{m_p\,m_{p+2}}$ ($\ell^2$ में कोशी--श्वार्ज़)। अतः
अनुपात $m_{p+1}/m_p$ $p$ में अनवरोही हैं; और चूँकि $R(x) = \frac{m_1}{m_0}$, $\frac{\langle Ax, Ax\rangle}
{\langle x, Ax\rangle} = \frac{m_2}{m_1}$ तथा $R(Ax) =
\frac{m_3}{m_2}$,
अतः शृंखला निकल आती है, जिसका प्रत्येक पद $\leq \mu_1$ है क्योंकि पदशः
$m_{p+1} \leq \mu_1m_p$। अभिसरण: यदि $a_1 > 0$ हो (शीर्ष अभिलक्षणिक मान का कुल भार $a_1$
लिखते हुए), तो योगों के प्रभावी अभिसरण से (अनुपात $< 1$)
\[
\mu_1 \geq R(A^kx) = \frac{m_{2k+1}}{m_{2k}} =
\mu_1\,\frac{a_1 + \sum_{\mu_n<\mu_1}(\mu_n/\mu_1)^{2k+1}
a_n}{a_1 + \sum_{\mu_n<\mu_1}(\mu_n/\mu_1)^{2k}a_n}
\longrightarrow \mu_1,
\]
अतः घात विधि प्रत्येक ऐसे आरंभिक सदिश के लिए अभिसरित होती है जो शीर्ष
अभिलक्षणिक समष्टि पर लांबिक न हो।

\textbf{12.} बिंदुशः $\abs{\langle x, (A - B)x\rangle}
\leq \vertiii{A - B}\,\norm x^2$, अतः प्रत्येक $x$ के लिए $R_A(x) \leq R_B(x) +
\vertiii{A - B}$। इसे
\cref{exo:b3:spectral:8} के उच्चिष्ठ--निम्निष्ठ सूत्र में डालने पर: $\mu_n(A)
\leq \mu_n(B) + \vertiii{A - B}$, और $A,
B$ में
सममित रूप से: अतः एक ही बार में सभी $n$ के लिए $\abs{\mu_n(A) - \mu_n(B)} \leq \vertiii{A - B}$।

\textbf{13.} $-w'' = x - x^2$ का समाकल $w = -\frac{x^3}6
+ \frac{x^4}{12} + cx$ है ($w(0) = 0$ के साथ), और $w(1) = 0$ $c = \frac1{12}$
दे देता है:
\[
w = \frac{x^4 - 2x^3 + x}{12} = \frac{x(1-x)(1 + x -
x^2)}{12} = Gu .
\]
तब $\norm u_2^2 = \int_0^1x^2(1-x)^2 = \frac1{30}$, और $\int_0^1x^3(1-x)^3 = B(4,4) = \frac1{140}$ के साथ:
\[
\langle u, Gu\rangle = \frac1{12}\Bigl(\frac1{30} +
\frac1{140}\Bigr) = \frac{17}{5040},
\qquad
R(u) = \frac{17/5040}{1/30} = \frac{17}{168} .
\]
अतः $\frac1{\pi^2} = \lambda_1 \geq \frac{17}{168}$, अर्थात् $\pi^2 \leq \frac{168}{17} = 9.8824$: $\pi \leq 3.14364 <
3.1437$ (वास्तविक मान $\pi^2 = 9.8696$)। एक बहुपद, एक
समाकल, एक अंक।

\textbf{14.} $u = x - x^2$ लिखिए, अतः $Gu =
\frac{u(1 + u)}{12}$ और $\int u^2 = \frac1{30}$, $\int u^3 = \frac1{140}$, $\int u^4 = B(5,5) =
\frac{4!\,4!}{9!} = \frac1{630}$ का उपयोग
करते हुए:
\[
\norm{Gu}_2^2 = \frac1{144}\int u^2(1+u)^2 =
\frac1{144}\Bigl(\frac1{30} + \frac2{140} +
\frac1{630}\Bigr) = \frac1{144}\cdot\frac{62}{1260} =
\frac{31}{90720} .
\]
अतः $\frac{\langle Gu, Gu\rangle}{\langle u, Gu\rangle} =
\frac{31/90720}{17/5040} = \frac{31}{306}$, और प्रश्न 11 से यह अब भी $\leq \mu_1 = \frac1{\pi^2}$ है: $\pi^2 \leq
\frac{306}{31} = 9.87097$, अर्थात् $\pi \leq 3.14181 <
3.1419$
--- चार अंक (और अगली पुनरावृत्ति लगभग आठ देती, क्योंकि त्रुटि प्रति
चरण $(\lambda_2/\lambda_1)^2 = \frac1{16}$ से सिकुड़ती है)।

\textbf{15.} कुल $(e_m \otimes e_n)(x,y) =
e_m(x)e_n(y)$ $L^2(\intcc01^2)$ का कोई \omterm{def:b3:hilbert:onb}{हिल्बर्ट आधार} है (टोनेली से
लांबिक-प्रसामान्यता; और \omterm{def:b3:complete:complete}{पूर्णता} \cref{exo:b3:spectral:5} की भाँति)। प्रश्न 7 से स्थिर
$x$ के लिए $g(x, \cdot) = \sum_n\lambda_ne_n(x)e_n$, अतः $e_m\otimes e_n$ पर $g$ का गुणांक
\[
\langle e_m\otimes e_n,\ g\rangle
= \int_0^1 e_m(x)\,\lambda_n e_n(x)\,\dd x
= \lambda_n\,\delta_{mn} .
\]
है। वर्ग में \omterm{thm:b3:hilbert:parseval}{पार्सेवाल}:
\[
\iint g^2 = \sum_{m,n}\abs{\langle e_m\otimes e_n,
g\rangle}^2 = \sum_n\lambda_n^2 .
\]

\textbf{16.} $g$ की सममिति से
\[
\iint g^2 = 2\iint_{y<x}y^2(1-x)^2
= 2\int_0^1(1-x)^2\,\frac{x^3}3\,\dd x
= \frac23\,B(4, 3) = \frac23\cdot\frac{3!\,2!}{6!}
= \frac23\cdot\frac1{60} = \frac1{90} .
\]

\textbf{17.} मिलाकर: $\sum_n\frac1{n^4\pi^4} =
\frac1{90}$, अर्थात् $\zeta(4) = \frac{\pi^4}{90}$। व्यापक रूप से $\operatorname{tr}(G^k) = \sum\lambda_n^k =
\frac{\zeta(2k)}{\pi^{2k}}$ $k$-घन
पर परिमेय-बहुपदीय अष्टि $g$ के गुणनफलों का कोई पुनरावृत्त समाकल
है: अर्थात् कोई परिमेय संख्या। अतः प्रत्येक $k$ के लिए $\zeta(2k) \in
\pi^{2k}\Q$। यह
मशीन केवल सम कोटियों तक पहुँचती है क्योंकि अभिलक्षणिक मान अपनी
\emph{घातों} से प्रवेश करते हैं --- $\sum\lambda_n^k$ --- और $\lambda_n =
\frac1{n^2\pi^2}$: संचिह्नों का
कोई संयोजन $\sum
n^{-3}$ उत्पन्न नहीं करता; और $\zeta(3)$ की अंकगणितीय प्रकृति
(अपेरी से अपरिमेय, अतिक्रांतता खुली) स्पेक्ट्रमी लेखा-जोखा से परे
है।

\textbf{18.} धनात्मक पदों के किसी योग का शीर्ष पद अधिक से अधिक वह
योग होता है: $\lambda_1^2 \leq \sum\lambda_n^2 =
\frac1{90}$, अतः $\frac1{\pi^2} \leq \frac1{\sqrt{90}}$ और $\pi \geq 90^{1/4} = 3.0801\ldots$ प्रश्न 14 के साथ: $3.080 < \pi < 3.1419$,
केवल तार-अंकगणित से। उच्चतर संचिह्न निचले परिबंध को ज्यामितीय रूप
से तीक्ष्ण करते हैं: $\lambda_1 \leq (\operatorname{tr}G^{2k})^{1/2k} =
\lambda_1\bigl(1 + \sum_{n\geq2}(\lambda_n/\lambda_1)^{2k}
\bigr)^{1/2k}$, और परजीवी गुणक $\bigl(\tfrac14\bigr)^{2k}\cdot\frac1{2k}$ की गति से मरता है
--- वही स्पेक्ट्रमी-रिक्ति क्रियाविधि जो घात विधि के अभिसरण की है
(प्रश्न 11), संचिह्न की ओर से देखी गई।

\textbf{19.} सभी $n$ के लिए $\abs{n^2\pi^2 - \nu} \geq \delta > 0$ (अनुक्रम $n^2\pi^2 \to \infty$ $\nu$ से एक अंतर
रखकर बचता है), और बड़े $n$ के लिए $\abs{n^2\pi^2 - \nu} \geq \frac{n^2\pi^2}2$। $L^2$ अभिसरण: गुणांक $\frac{c_n}{n^2\pi^2 - \nu}$
वर्ग-योगनीय हैं ($\frac{\abs{c_n}}\delta$ का प्रभुत्व)। एकसमान अभिसरण: पुच्छ उच्चतम-मानक
$\sqrt2\sum_{n>N}
\frac{\abs{c_n}}{\abs{n^2\pi^2 - \nu}} \leq
\frac{2\sqrt2}{\pi^2}\bigl(\sum\abs{c_n}^2\bigr)^{1/2}
\bigl(\sum_{n>N}n^{-4}\bigr)^{1/2} \to 0$ से परिबद्ध हैं (कोशी--श्वार्ज़)। सत्यापन: $Gf + \nu Gu$ का $e_n$-गुणांक
\[
\lambda_nc_n + \frac{\nu\lambda_nc_n}{n^2\pi^2 - \nu}
= \frac{c_n}{n^2\pi^2}\Bigl(1 + \frac{\nu}{n^2\pi^2 -
\nu}\Bigr) = \frac{c_n}{n^2\pi^2 - \nu} :
\]
है, जो ठीक $u$ के गुणांक हैं, अतः $u = Gf + \nu Gu$; और अद्वितीयता इसलिए कि
हलों का कोई अंतर $v$ $v = \nu
Gv$ संतुष्ट करता है, अर्थात् सभी $n$ के
लिए $\langle e_n, v\rangle(n^2\pi^2 - \nu) = 0$: $v = 0$।

\textbf{20.} प्रश्न 19 की भाँति समीकरण $u = Gf + \nu
Gu$ गुणांक समीकरणों $(n^2\pi^2 - \nu)\,\langle e_n, u\rangle = c_n$,
$n \geq
1$ के कुल के तुल्य है। $n = m$ के लिए बायाँ पक्ष $0$ है: अतः
हल-योग्यता $c_m = 0$ अनिवार्य कर देती है, और तब $\langle e_m,
u\rangle$ स्वतंत्र रहता है
जबकि शेष सभी गुणांक निर्धारित होते हैं: अतः हल रेखा $u_0 + \R e_m$ बनाते हैं।
अनुनाद: अभिलक्षणिक विधा पर घटक रखने वाला कोई बल उसमें बिना किसी
परिबंध के ऊर्जा भरता जाता है --- अपनी ही आवृत्ति पर धकेला गया झूला।

\textbf{21.} प्रश्न 19 के सूत्र से $\norm{R_\nu f}_2^2
= \sum\frac{\abs{c_n}^2}{(n^2\pi^2 - \nu)^2} \leq
\frac{\norm f^2}{(\pi^2 - \nu)^2}$ ($\nu < \pi^2$ के लिए निकटतम
अभिलक्षणिक मान $\pi^2$ है), जिसमें बराबरी $f = e_1$ पर निकट आती है: अतः
\omterm{def:b3:banach:operator}{संकारक मानक} $\frac1{\pi^2 - \nu}$। संहतता: $R_\nu$ अपने परिमित-कोटि कर्तनों की मानक-सीमा
है (पुच्छ गुणांक $\frac1{n^2\pi^2 - \nu} \to 0$); और \omterm{def:b3:spectral:selfadjoint}{स्वसंलग्नता} तथा धनात्मकता विकर्ण रूप से
पढ़ ली जाती हैं (सभी गुणांक $\frac1{n^2\pi^2 - \nu} > 0$)। $R_\nu$ के अभिलक्षणिक मान $\frac1{n^2\pi^2 - \nu}$
हैं: अतः भाग I से VI तक का विश्लेषण अक्षरशः फिर से चल पड़ता है।

\textbf{22.} शब्दकोश: \emph{अभिलक्षणिक मान} $\lambda_n =
\frac1{n^2\pi^2}$ $\leftrightarrow$ $n$-वें
संनादी की वर्ग आवृत्ति $n^2\pi^2$; \emph{संचिह्न} $\sum
\lambda_n = \frac16$ $\leftrightarrow$ $\zeta(2) =
\frac{\pi^2}6$;
\emph{हिल्बर्ट--श्मिट मानक} $\iint g^2 =
\frac1{90}$ $\leftrightarrow$ $\zeta(4) = \frac{\pi^4}{90}$; \emph{\omterm{thm:b3:spectral:fredholm}{फ्रेडहोम विकल्प}}
$\leftrightarrow$ बलित तार का अनुनाद; \emph{निम्निष्ठ--उच्चिष्ठ} $\leftrightarrow$ विचरणात्मक
आकलन, एक ही बहुपद से $\pi < 3.1437$ तक। प्रत्येक जोड़े के पीछे वही वस्तु है:
एक ही \omterm{def:b3:spectral:compact}{संहत} \omterm{def:b3:spectral:selfadjoint}{स्वसंलग्न संकारक}, एक बार विकर्णीकृत और पाँच तरह से उपयोग
किया गया।

\textbf{23.} सम-विषमता पर विभाजित करके और सम भाग में $n = 2m$ प्रतिस्थापित
करके
\[
\zeta(6) = \sum_{n\text{ विषम}}\frac1{n^6} +
\sum_{m\geq1}\frac1{(2m)^6}
= \frac{\pi^6}{960} + \frac{\zeta(6)}{64},
\]
अतः $\frac{63}{64}\zeta(6) = \frac{\pi^6}{960}$ और $\zeta(6)
= \frac{64\,\pi^6}{63\cdot960} = \frac{\pi^6}{945}$। संचिह्न वाला मार्ग $\operatorname{tr}G^3 =
\sum\lambda_n^3 = \zeta(6)/\pi^6$ को पुनरावृत्त अष्टि
$g_2(x,y) = \int_0^1g(x,z)g(z,y)\dd z$ के साथ $\iint g\,g_2$ के रूप में परिकलित करता --- अर्थात् खंडशः बहुपदों
के तीन समाकलन; जबकि $x(1-x)$ पर \omterm{thm:b3:hilbert:parseval}{पार्सेवाल} को केवल एक चाहिए था।

\textbf{24.} (क) विकर्णीकरण कीजिए: $v = \sum_na_nu_n$ (और साथ में कोई संभावित
अष्टि घटक, जिस पर $\langle v, Av\rangle$ कुछ नहीं पाता और $\norm v^2$ बढ़ता है, अतः किसी
अधिकतमकारी में वह नहीं होता)। तब $\langle v, Av\rangle = \sum\mu_na_n^2 \leq
\mu_1\sum a_n^2$, जिसमें बराबरी तभी और केवल
तभी जब $\mu_n < \mu_1$ होने पर $a_n = 0$: अर्थात् कोई अधिकतमकारी $\mu_1$-अभिलक्षणिक
समष्टि में होता है। (ख) किसी भी $u$ के लिए
\[
\langle\abs u, A\abs u\rangle - \langle u, Au\rangle
= \iint k(x,y)\,\bigl(\abs{u(x)}\abs{u(y)} -
u(x)u(y)\bigr)\dd x\,\dd y \;\geq\; 0,
\]
जिसका समाकल्य बिंदुशः ऋणेतर है। यदि $P = \{u >
0\}$ और $N = \{u < 0\}$ दोनों का \omterm{def:b3:measure:measure}{माप}
धनात्मक हो, तो $P\times N$ पर समाकल्य धनात्मक \omterm{def:b3:measure:measure}{माप} के किसी समुच्चय पर $2k\abs{u(x)}\abs{u(y)} >
0$
के बराबर है: अतः कड़ी असमिका। कोई $\mu_1$-अभिलक्षणिक फलन $u$ रेली
भागफल को अधिकतम कर देता है, और $\abs u$ का मानक वही है, अतः कड़ाई $R(\abs u) > \mu_1$
प्रस्तुत कर देती --- असंभव; इसलिए $u$ का चिह्न लगभग सर्वत्र अचर
है, कहिए $u \geq 0$। तब प्रत्येक $x \in
\intoo01$ के लिए $u(x) = \mu_1^{-1}(Au)(x) =
\mu_1^{-1}\int k(x,y)u(y)\dd y > 0$ ($k(x,\cdot) > 0$ और $u \neq 0$)।
(ग) यदि अभिलक्षणिक समष्टि की विमा $\geq 2$ होती, तो उसमें दो \emph{लांबिक}
अभिलक्षणिक फलन $u, v$ होते, जिनमें से प्रत्येक (ख) से अचर चिह्न वाला
और \omterm{def:b3:topology:interior}{अंतःभाग} में अलुप्त होता; पर तब $\abs{\langle u, v\rangle} = \int\abs u\,\abs v > 0$ --- विरोधाभास। तार पर: विवृत
वर्ग पर $k = g > 0$, $\mu_1 = \lambda_1 = \frac1{\pi^2}$ वस्तुतः सरल है, $e_1 = \sqrt2\sin(\pi x)$ $\intoo01$ पर धनात्मक है; और
धनात्मक $e_1$ पर लांबिक प्रत्येक $e_n = \sqrt2\sin(n\pi x)$, $n \geq 2$, को उसके सापेक्ष शून्य
पर समाकलित होना पड़ता है, अतः वह चिह्न बदलता है --- जैसा उसके $n - 1$
आंतरिक शून्य $\frac kn$ पुष्टि करते हैं।

\textbf{25.} चूँकि $n^2\pi^2 \to \infty$, अतः निम्निष्ठ $d =
\min_n\abs{n^2\pi^2 - \nu}$ किसी विधा $m$ पर
प्राप्त होता है, और $d > 0$ क्योंकि $\nu$ स्पेक्ट्रम से बचता है। प्रश्न
19 का विकर्ण सूत्र
\[
\norm{R_\nu f}_2^2 =
\sum_n\frac{\abs{c_n}^2}{(n^2\pi^2 - \nu)^2}
\leq \frac1{d^2}\,\norm f_2^2,
\]
देता है, जिसमें $f = e_m$ के लिए बराबरी: $\vertiii{R_\nu} = \frac1d$, अर्थात् $\nu$ की स्पेक्ट्रम
से दूरी का व्युत्क्रम --- वही व्यापक विलायक सिद्धांत, यहाँ स्पष्ट
निर्देशांकों में। \omterm{def:b3:spectral:selfadjoint}{स्वसंलग्नता} वास्तविक विकर्ण गुणांकों से पढ़ ली जाती
है; और संहतता प्रश्न 21 की भाँति निकलती है (गुणांक $0$ की ओर जाते
हैं, अतः परिमित-कोटि कर्तन मानक में अभिसरित होते हैं)। अनुनाद की
क़ीमत: $f = e_1$ और $\nu =
(1-\varepsilon)\pi^2$ के लिए सूत्र $u =
\frac{c_1}{\pi^2 - \nu}e_1 = \frac1{\varepsilon\pi^2}e_1$ देता है, जबकि स्थैतिक अनुक्रिया
$Ge_1 = \frac1{\pi^2}e_1$ है: अतः आवर्धन $\frac1\varepsilon$। $\varepsilon =
10^{-2}$ पर अनुक्रिया स्थैतिक की $100$ गुना
है --- और वह $\varepsilon \to 0$ पर अपसारित हो जाती है, जो प्रश्न 20 का विकल्प
परिबद्ध ओर से देखा गया है: बल की आवृत्ति किसी प्राकृतिक आवृत्ति के
जितनी निकट होगी, प्रतिलोम उतना ही कम परिबद्ध।
\end{solution}
